% =====================================================================
% Chapter 7 -- Discussion
%
% Outline: outlines/ch7_discussion.md (revision 2, approved 2026-10-24; Sec. 7.1 revision 4, 2026-11-06; Sec. 7.3 rewrite outline, approved 2026-11-10).
% Review fixes 2026-11-10 (author request "fix all issues except 3"): backup backups/discussion_2026-11-10_before_review_fixes.tex
% Interview answers: thesis_notes.md §1 (2026-10-24).
% No new results: every number is taken from Chapter 6, which names its source in results/.
% Every paragraph marked CLAUDE-DRAFT is to be rewritten by the author.
% Review fixes 2026-11-06 (author request "fix all"): backup backups/discussion_2026-11-06_before_review_fixes.tex
% Language and cutting pass 2026-11-06: backup backups/discussion_2026-11-06_before_language_pass.tex
% =====================================================================

\chapter{Discussion}
\label{chap:discussion}

% CLAUDE-DRAFT
This chapter interprets the results of \cref{chap:experiments} across the experiments.
\Cref{sec:disc-rq} answers the research question, \cref{sec:disc-limitations} states the limitations and threats to validity, and \cref{sec:disc-generalisation} covers generalisation and extensions.

% ---------------------------------------------------------------------
\section{Answer to the Research Question}
\label{sec:disc-rq}
% Redrafted 2026-11-06 (round 2) from the author's interview answers Q1-Q18 (thesis_notes.md §1, rows 'Sec. 7.1 round 2'); outline revision 4 approved 2026-11-06. Organised by theme. Backup: backups/discussion_2026-11-06_before_7.1_round2.tex

% CLAUDE-DRAFT
The research question (\cref{sec:intro-rq}) asks how continual-learning strategies trade adaptation against forgetting.
Among the strategies that use one parameter set and train only on the data of the current task, the more a strategy adapted with the full training split, the lower its Final Performance (\cref{tab:ch6_main}).
In the original task order, no run of these strategies exceeded a Final Performance of 0.233 at any rank or number of training volumes, although their Adaptation reached up to 0.781 (\cref{tab:ch6_few_shot,tab:ch6_rank}).
After the last task, these strategies segmented little beyond the one or two structures trained last (\cref{fig:ch6_main_matrices}).
Shared LoRA also kept femoral cartilage, the one structure that the pretrained model already segmented.
% source: eval_matrix.csv, last column: seq. FT (a9f4ca56) and distillation-based (65011743) only lateral meniscus (0.74, 0.80), all others <= 0.02; prompt tuning (7bcbb6cc) lateral meniscus 0.53, medial meniscus 0.33, others <= 0.04; shared LoRA r = 1024 (90e8b4e6) femoral 0.55, medial meniscus 0.36, lateral meniscus 0.40, others <= 0.04; pretrained: femoral 0.68, all other knee structures 0.00

% CLAUDE-DRAFT
In the original task order, only shared LoRA with replay and task-specific LoRA ended above a Final Performance of 0.233.
Replay raised the Final Performance of shared LoRA without lowering its Adaptation, but it requires storing patient data from earlier tasks and more training time per epoch.
Task-specific LoRA keeps a separate adapter for each task and therefore requires the task identity at inference.
Both strategies adapted less than sequential fine-tuning.
These results agree with the Beyond Forgetting benchmark on 2D medical segmentation, in which replay gave the best balance between retaining earlier tasks and learning new ones, and parameter isolation avoided forgetting at the cost of a growing model~\cite[Sec.~5]{wang2026beyondforgetting}.
In MedPEFT-CL, which adapts a 2D text-prompted model with LoRA, replay likewise reduced forgetting more than any other component~\cite[Table~8]{gao2026medpeftcl}, and in Lifelong nnU-Net, which trains a 3D nnU-Net on sequences of \gls{mri} datasets, it forgot the least~\cite{gonzalez2023lifelong}.
% source: paper_notes/cards/NEW_beyondforgetting2026.md l. 17 (Abstract; Sec. 5, p. 10: replay best stability-plasticity balance; parameter isolation removes forgetting at the cost of growing model size); method.tex l. 274 (task-specific LoRA evaluates the control task without an adapter); Lifelong nnU-Net: paper_notes/gonzalez2023lifelong.md l. 46, 50 (all 3D, nnU-Net 3D full-resolution), related_work.tex l. 205 (rehearsal least forgetting)

% CLAUDE-DRAFT
Neither retaining the earlier tasks nor freezing the pretrained weights kept the pretrained ability.
In the comparison of the strategies, shared LoRA with replay forgot the least of those with one parameter set but still lost part of its performance on the control task (\cref{sec:exp-main}).
The prompt-tuning strategy lost most of it, although its pretrained weights stay frozen.
Only task-specific LoRA kept the control task, because this strategy evaluates it with the pretrained model without an adapter.
% source: ch6_main.tex (Adaptation / Final Performance: shared LoRA 0.471 / 0.233, prompt tuning 0.580 / 0.155, distillation-based 0.747 / 0.135, seq. FT 0.781 / 0.130; replay 0.472 / 0.450; task-specific 0.485, highest); cl_metrics.csv, all reported runs in the original order with one parameter set and current-task data only (seq. FT full, k = 5/10/20, equal updates; distillation-based; prompt tuning; shared LoRA r = 4...2048, k = 5/10/20, equal updates): max Final Performance 0.233 (shared LoRA r = 2048: 0.2330, r = 1024: 0.2326); runs above 0.233 in the original order: task-specific LoRA r = 4 (0.485), r = 8 (0.404), equal updates (0.429), replay (0.450); Sec. 6.7 and Sec. 5.3 (replay epoch time, stored data; task identity); paper_notes/NEW_medpeftcl2025.md (Table 8, p. 8: cumulative ablation, FR 20.32 -> 14.01 with EWC, 14.01 -> 6.56 with replay memory, -> 3.38 Fisher guidance, -> 1.91 bidirectional feedback; replay = largest single drop in FR)

% CLAUDE-DRAFT
Under few-shot conditions, with 5 to 20 volumes and ten epochs per task, sequential fine-tuning still learned every task, whereas the LoRA strategies hardly adapted (\cref{sec:exp-few-shot}).
With 5 volumes, sequential fine-tuning also forgot less of the earlier tasks than with the full training split and kept the control task.
At $k = 5$, the equal-updates runs point to the number of updates rather than the number of volumes as the cause of the lower forgetting, the kept control task and the weak adaptation of the LoRA strategies.
The LoRA strategies therefore needed more updates rather than fewer to learn from few volumes.
In the few-shot study of Silva-Rodríguez et al., which used a different model and different data, LoRA learned new structures from 10 labelled volumes when trained for up to 200 epochs~\cite[Sec.~4.3.2]{silvarodriguez2023fseft}.
% source: ch6_few_shot.tex (seq. FT 5 (50): Final Forgetting 0.382 vs 0.782; equal updates: LoRA Adaptation 0.429 / 0.438, seq. FT Final Forgetting 0.714); papers/t2_few_shot/NEW_silvarodriguez2023fseft.pdf, Sec. 4.3.2 (up to 200 epochs), checked 2026-11-06; moved here from Sec. 7.2 on 2026-11-10

% CLAUDE-DRAFT
Apart from the strategy itself, the task order and the rank of shared LoRA also shaped the trade-off.
The task order changed the results of shared LoRA but not the ranking of the three strategies that we ran in both orders (\cref{sec:exp-task-order}).
In the reverse order, shared LoRA ended with a Final Performance of 0.345, the only run with one parameter set and current-task data above 0.233 (\cref{tab:ch6_task_order}).
Over ten orders of 2D prostate segmentation tasks, the Beyond Forgetting benchmark likewise found that the task order affects some strategies more than others~\cite[Sec.~4.4]{wang2026beyondforgetting}.
A ranking obtained from one or two orders may therefore not hold in others.
% source: ch6_task_order.tex (shared LoRA original / reverse: Final Forgetting 0.289 / 0.124, Adaptation 0.471 / 0.439); paper_notes/cards/NEW_beyondforgetting2026.md (Sec. 4.4, Fig. 2, p. 9-10: Domain-CL, ten random orders; EWC and LwF vary most, GSS least)

% CLAUDE-DRAFT
Within shared LoRA, higher ranks forgot less at a lower Adaptation, with the scaling factor $\gamma_r$ falling as the rank rose (\cref{sec:exp-rank}).
The fall in Adaptation contrasts with rsLoRA, whose scaling factor we use and which improved with increasing rank when fine-tuning a large language model~\cite[Sec.~4]{kalajdzievski2023rankstabilizationscalingfactor}.
We did not test which of the differences in model, task and training configuration explains the opposite trend.

% CLAUDE-DRAFT
The distillation-based strategy was less effective in our setting than ZSCL in its own evaluation.
In ZSCL, adding distillation to sequential fine-tuning raised Transfer, the performance on tasks before they are trained, from 44.6 to 58.9\,\%~\cite[Table~2]{zheng2023preventingzeroshottransferdegradation}, whereas our distillation-based strategy reached 0.002, against 0.001 for sequential fine-tuning (\cref{sec:exp-main}).
Of the differences between the two methods (\cref{subsec:method-distillation}), distillation on current-task data is unlikely to explain the gap on its own, because in ZSCL's ablation it kept Transfer at the level of a reference dataset (\cref{sec:cl-vlm}).
% source: paper_notes/zscl.md l. 33 (Table 1b: current task 56.7 vs ImageNet 56.8 Transfer), l. 41 (Table 2: 44.6 -> 58.9), l. 46 (Transfer = tasks not yet trained); related_work.tex l. 116; method.tex l. 374-376 (differences); ch6_main.tex (Transfer: distillation-based 0.002, seq. FT 0.001)
The untuned weight $\lambda = 0.01$, the missing weight-consolidation loss and the move from classifying images to segmenting 3D volumes remain as possible causes, and we did not test which one explains the gap.

% CLAUDE-DRAFT
All of these comparisons rest on the evaluation protocol (\cref{sec:method-setting-protocol}), which measures the loss of pretrained ability separately from the forgetting of earlier tasks.
The two did not always move together.
From $r = 1024$ to $r = 2048$, the Final Forgetting of shared LoRA hardly changed, whereas its control-task drop fell from 0.237 to 0.091 (\cref{sec:exp-rank}).
Conversely, sequential fine-tuning with 20 volumes per task forgot less of the earlier tasks than with the full training split but lost the control task just as completely (\cref{sec:exp-few-shot}).
Neither metric can therefore stand in for the other.
Our control task follows the control set of SPU~\cite{Zhang_2024_CVPR}, and comparable measures exist for segmentation models prompted with points or boxes~\cite{wang2025casam,wang2024sammed3dmoe}, but none of the works closest to ours uses one (\cref{sec:sota-summary}).
The pretrained-model reference showed that the LoRA strategies with 5 and 10 volumes per task did not learn at all, as their Adaptation hardly differed from that of the pretrained model (\cref{sec:exp-few-shot}).
Without this reference, their Adaptation would read as weak learning.
% source: thesis/tables/ch6_rank.tex (shared LoRA r = 1024 / 2048: Final Forgetting 0.289 / 0.284, control-task drop 0.237 / 0.091), ch6_few_shot.tex (seq. FT 20 (200) / 108 (1,080): Final Forgetting 0.573 / 0.782, control-task drop 0.979 / 0.979; LoRA 5 (50), 10 (100)); related_work.tex l. 271 (closest works), introduction.tex l. 102 (none measures pretrained ability)

% ---------------------------------------------------------------------
\section{Limitations and Threats to Validity}
\label{sec:disc-limitations}

% CLAUDE-DRAFT
Three limitations affect the conclusions most: the shared training configuration, the single run per configuration and the narrow scope of the study.
The remaining limitations concern the strategies we did not evaluate, the evaluation protocol, the setup of the runs and the data.

% CLAUDE-DRAFT
First, all strategies were trained with one configuration, and neither the learning rate nor the number of epochs was tuned per strategy, so our comparison measures the strategies under one budget, not each at its best setting (\cref{sec:method-adaptation}).
Three experiments also change a second factor together with the one they study: the few-shot settings reduce the number of updates together with the number of volumes, the shared-LoRA rank experiment changes the scaling factor $\gamma_r$ together with the rank, and the 50-epoch patellar run decays the learning rate more slowly, together with the larger number of updates (\cref{sec:exp-few-shot,sec:exp-rank,sec:exp-patellar}).
In these experiments, a difference cannot be attributed to the studied factor alone, except for volumes and updates at $k = 5$, which the equal-updates runs separate for three strategies.

% CLAUDE-DRAFT
Within the ten epochs per task used for all strategies, shared LoRA with $r = 1024$ learned little on its first task in both task orders, and task-specific LoRA learned nothing on patellar cartilage at any rank tested, as in the few-shot settings (\cref{sec:exp-task-order,sec:exp-patellar,sec:disc-rq}).
We see two possible explanations, but none of our runs tests either of them.
The first is the zero initialisation of $\mat{B}$, which leaves the prediction of the model unchanged at the start of training (\cref{subsec:method-lora}).
The second is a training setup that was not tuned for LoRA: the adapters were trained with stochastic gradient descent, whereas Hu et al.\ motivate the scaling of LoRA for training with Adam~\cite[Sec.~4.1]{hu2022lora}, and the scaling factor $\gamma_r$ of shared LoRA falls as the rank rises.
The low Adaptation of the LoRA strategies may therefore reflect our training setup as much as the strategies themselves.\todo{defense: zero-initialisation and training-setup explanations untested; tests in \cref{sec:disc-generalisation}}
% source: Ch. 6.6 l. 450 (shared LoRA r = 4, first task: loss 0.38 after ten epochs; r = 1024: 0.69); tab:app-strategy-params (gamma_r 16 at r = 4, 1 at r = 1024); tab:app-params (SGD, Nesterov, eta_0 = 1e-4); paper_notes/lora_hu.md l. 16 (Sec. 4.1: with Adam, tuning alpha ~ tuning the learning rate), l. 31 (SGD instead of Adam)

% CLAUDE-DRAFT
Second, each configuration was trained once, on one random split of each dataset and, in the few-shot settings, on one draw of training volumes, so our bootstrap intervals cover only the choice of test volumes (\cref{subsec:exp-metrics}).
The two task-specific LoRA runs of the task-order experiment differ only in the initialisation of the adapters, yet the bootstrap interval of their difference in Adaptation excludes zero (\cref{sec:exp-task-order}).
% source: thesis/tables/cl_differences.csv (adaptation, TIL-LoRA r = 4 original - reverse: 0.006 [0.004, 0.008]); Ch. 6 l. 413-419
Small differences between single runs can therefore come from training randomness alone, even when their interval excludes zero.
% source: thesis/tables/ch6_rank.tex (shared LoRA Final Forgetting 0.546 -> 0.284, control-task drop 0.978 -> 0.091), ch6_task_order.tex (shared LoRA Final Forgetting 0.289 / 0.124)
The effects of the rank and the task order on shared LoRA are many times larger than this difference, but we did not repeat any configuration of shared LoRA, so its variation between runs is unknown.
Our bootstrap intervals are also not adjusted for multiple comparisons.

% CLAUDE-DRAFT
Third, the scope of the study is narrow.
We adapted one model, VoxTell, to six structures of one knee \gls{mri} dataset in two task orders, without data augmentation and with one fixed text prompt per structure for both training and evaluation (\cref{sec:data-skmtea}).
Our evaluation therefore cannot show whether the adapted models remain as robust to rephrased text prompts as the pretrained VoxTell (\cref{sec:text-promptable-seg}).
The distillation-based and the prompt-tuning strategies were run only with the full training split and in the original task order, so the few-shot part of the research question is answered only for sequential fine-tuning, task-specific LoRA and shared LoRA (\cref{tab:ch6_experiments}).
Shared LoRA with replay was likewise run only with the full training split, in the original task order and at $r = 1024$.
The effect of replay in other configurations is unknown, because the task order and the rank strongly changed how shared LoRA balanced adaptation and retention (\cref{sec:exp-rank,sec:exp-task-order}).
The same uncertainty applies to the smaller control-task drop with replay, which arose in the last task and whose cause we did not test (\cref{sec:exp-main}).

% CLAUDE-DRAFT
Our results also cover only the strategies we evaluated, in the form we implemented them.
The distillation-based and prompt-tuning strategies are adaptations of ZSCL and CPE-CLIP (\cref{subsec:method-distillation,subsec:method-prompt}), so their results do not describe the original methods.
Our distillation teacher is the pretrained model, as in ZSCL, and we did not test other teachers, such as the model of the previous task used by PRIMED (\cref{sec:cl-vlm}).
We also did not evaluate sharing adapters between tasks with similar text prompts, as MedCRP-CL does by clustering the tasks by the embeddings of their text prompts~\cite{gao2026medcrpcl}, or the other methods named in \cref{sec:method-strategies}.

% CLAUDE-DRAFT
The evaluation protocol has two further limits.
The control task measures the loss of pretrained ability with a single large, high-contrast structure, the lung in \gls{ct}, so it cannot show how the strategies affect the many other structures and modalities that VoxTell was trained on.
With only four test volumes, its \gls{dsc} is also less certain than that of the SKM-TEA tasks, which are evaluated on 24 volumes.
The origin of its reference labels is not documented (\cref{sec:data-aeropath}).
Because AeroPath is among the training sources of checkpoint v1.1, the control task also measures the loss on data the model was trained on, not on unseen data (\cref{sec:data-aeropath}).
The second limit is the joint reference: it covers all six tasks only, so the intransigence measure of Chaudhry et al.\ cannot be computed (\cref{sec:method-setting-protocol}).

% CLAUDE-DRAFT
Three further threats to internal validity come from how the runs were set up and evaluated.
The ranks of the two LoRA strategies in the comparison were chosen from the test results of the rank experiment (\cref{subsec:exp-config}), so the results favour these two configurations.\todo{defense: main-comparison ranks chosen on test-split results (optimistic bias); repeated from Sec. 6.1.2}
The two LoRA strategies were also run with different values of $\alpha$ and of the dropout rate (\cref{subsec:method-lora}), so a difference between them cannot be attributed to the sharing of the adapter alone.\todo{defense: alpha and dropout differ between the two LoRA variants (moved from Sec. 6.1.2)}
The choice of the evaluated checkpoint can also reverse close comparisons.
At the last epoch, the distillation-based strategy forgot 0.048 less than sequential fine-tuning, but at the checkpoint with the lowest validation loss it forgot 0.021 more $[0.012, 0.029]$ (\cref{tab:ch6_main,tab:ch6_components}).
% source: thesis/tables/cl_differences.csv (exp_20261022_171452_65011743:best_validation - exp_20261022_171452_a9f4ca56:best_validation, final_forgetting +0.021 [0.012, 0.029]); ch6_main.tex, ch6_components.tex

% CLAUDE-DRAFT
Finally, the data and the checkpoint limit how far our numbers carry over to other evaluations.
We used our own random split of SKM-TEA instead of the official one, so our numbers cannot be compared with published SKM-TEA results.
The SKM-TEA files store no orientation, so the volumes reach the model in their stored axis order, which may differ from that of the knee volumes in VoxTell's training data and may lower the results of the pretrained model (\cref{sec:data-preprocessing}).
The model also sees only the first \gls{qdess} echo, although the annotators used both echoes to separate neighbouring cartilage and meniscus voxels~\cite[App.~A.4]{desai2021skmtea}.
On our test volumes, the pretrained model segmented femoral cartilage with a \gls{dsc} of 0.68 but none of the other five structures (\cref{sec:exp-main}), whereas the VoxTell authors report 46.45\,\% for femoral cartilage and at most 14.78\,\% for the others with checkpoint v1.0~\cites[Suppl.~Table~9]{rokuss2025voxtellfreetextpromptableuniversal}{rokuss2026voxtellmodelcard}.
% source: Pretrained column of eval_matrix.csv (e.g. exp_20260926_095539_90e8b4e6); paper_notes/voxtell.md (Table 9)
We use checkpoint v1.1, which differs from v1.0 in its training data.
Its added sources include ``Stanford Knee''~\cite{rokuss2026voxtellissue15}, the name the VoxTell paper uses for SKM-TEA, so our test volumes may have been part of its training data, although the pretrained model did not segment five of the six structures.\todo{defense: v1.1 lists Stanford Knee (= SKM-TEA in the VoxTell paper) among added training sources; volumes and labels unknown}
% source: github.com/MIC-DKFZ/VoxTell/issues/15 ('Stanford Knee' in 'Datasets added in v1.1'); VoxTell paper main text and Suppl. ('Stanford Knee dataset [50]', [50] = SKM-TEA); huggingface.co/mrokuss/VoxTell (v1.0 'used for the experiments in the paper'), checked 2026-11-07

% ---------------------------------------------------------------------
\section{Generalisation and Extensions}
\label{sec:disc-generalisation}

% Rewritten 2026-11-10 (outline approved by the author: outlines/ch7_discussion.md, 'Sec. 7.3 rewrite outline'). Backup: backups/discussion_2026-11-10_before_7.3_rewrite.tex

% CLAUDE-DRAFT
The numerical results and the ranking of the strategies apply only to the setting studied here (\cref{sec:disc-limitations}).
The evaluation protocol (\cref{sec:method-setting-protocol}) depends less on this setting.
It requires only a task sequence, a held-out control task that the pretrained model segments well, and the pretrained model and joint fine-tuning as references.
It can therefore be applied to other free-text promptable models, where it would show which of our observations depend on VoxTell, and to other task sequences.

% CLAUDE-DRAFT
The contrast between sequential fine-tuning and task-specific LoRA (\cref{sec:disc-rq}) has also been reported in another setting and may therefore hold more generally.
With a 3D brain-\gls{mri} model, two tasks and 16 or 32 labelled samples per task, Chen et al.\ found that sequential fine-tuning forgot most of the first task, whereas task-specific LoRA modules kept it~\cite{chen2026fewshotcontinuallearning3d}.
Because such modules keep earlier tasks by construction, this agreement concerns mainly the forgetting of sequential fine-tuning.

% CLAUDE-DRAFT
The main extension is a strategy that keeps the pretrained ability without being given the task identity at inference.
In the comparison of the strategies, every strategy with one parameter set lost part or all of the control task, whereas task-specific LoRA kept it but needs the task identity to select its adapter (\cref{sec:disc-rq}).
On this single structure, adaptation therefore cost the model part of its universality.
With one fixed text prompt per structure, the task identity can be derived from the text prompt, but a rephrased text prompt or a new structure has no adapter of its own.
One remedy is to select the parameters from the input itself, as MoE-Adapters does~\cite{yu2024boostingcontinuallearningvisionlanguage}, although MedCRP-CL outperformed it on 2D medical segmentation~\cite{gao2026medcrpcl}.
Another remedy is to select the adapter of task-specific LoRA by the similarity of the text prompt to those of the trained tasks, in the way MedCRP-CL groups tasks (\cref{sec:disc-limitations}).
Such a selection could also assign an adapter to a rephrased text prompt.
A further candidate is nested learning, which represents a model as a set of nested optimisation problems~\cite{behrouz2025nestedlearning}.
Its authors propose a memory of blocks that are updated at different frequencies, so that knowledge lost in one block may remain in others, and evaluate it on class-incremental text classification, in which the task identity is not given at inference.
However, nested learning was developed and evaluated with language models, and whether it carries over to 3D segmentation is untested.
% sources: related_work.tex l. 138-139 (MoE-Adapters, autoencoder selection), l. 245-246 (MedCRP-CL 73.3 vs MoE-Adapters 65.3, clustering by CLIP embeddings of text prompts); paper_notes/nested_learning.md

% CLAUDE-DRAFT
Shared LoRA with replay could be extended in two ways.
It was run only in the original task order and at $r = 1024$ (\cref{sec:disc-limitations}), so runs in the reverse task order and at other ranks would show whether its effect holds in other configurations.
It also draws the stored samples at random (\cref{subsec:method-replay}), whereas Bera et al.\ and MedPEFT-CL select them by criteria such as training difficulty (\cref{subsec:cl-medseg-net,subsec:cl-medseg-text}), and such a selection could be compared with random selection.

% CLAUDE-DRAFT
Further extensions concern the training and the number of runs.
The two explanations for the slow learning of the LoRA adapters (\cref{sec:disc-limitations}) can be tested by training the first task of shared LoRA at $r = 1024$ for more epochs, by initialising the adapters of task-specific LoRA from an adapter trained on another structure instead of from zero, and by varying the learning rate or $\alpha$ at a fixed high rank.
Varying $\alpha$ at a fixed rank would also separate the effect of the rank from that of the scaling factor $\gamma_r$.
The other limitations in \cref{sec:disc-limitations} point to direct extensions: equal-updates runs with 10 and 20 volumes per task, few-shot runs of the distillation-based and prompt-tuning strategies, tuning the learning rate, the number of epochs and the LoRA ranks on the validation split, and repeated runs with several seeds and draws of training volumes.

% CLAUDE-DRAFT
Other extensions concern the evaluation.
More task orders would show whether the forgetting and the control-task drop of shared LoRA in the last task of the original order come from training on lateral meniscus or from the tasks trained before it (\cref{sec:exp-task-order}).
Because all six tasks use the same images (\cref{sec:method-setting-protocol}), task sequences from other datasets, body regions and modalities would show whether the results hold when the images change between tasks.
Rephrased text prompts, for instance in the clinical styles of SEER~\cite{zhang2026seer}, would show whether adaptation with one text prompt per structure changes how the model segments that structure when it is named differently.
A control task with several structures and modalities would measure the loss of pretrained ability more broadly than our control task, which contains only the lung.
